% Nouveau EGA V (cote 25): archive cote 25, site page 118 (PDF page 119).
% typescript. Draft compared: none.
% AI-transcribed by Claude Opus 5.5 (claude-opus-5-5, adaptive thinking, effort high), Message Batches API, batch ega25b (msgbatch_01A7n5YuXubUPNK1zsepi8RG), 2026-11-13; checked against both supplied scan parts; see CHECK_DECISIONS.tsv.
\scanpage{25/118}
\mspage{\hand{21}}

\placedmargin{$|$}\placedmargin{$\wr$}ce qui entrainera la conclusion voulue. Soit $\underline{J}$ l'Idéal qui définit
D , il suffira évidemment de prouver que $\underline{J}\hand{\to} i_{\hand{*}}(\underline{J}\hand{|}X-Z")$ est un isomor
phisme ($i:X-Z"\hand{\to} X$ désignant l'immersion canonique), alors l'homomorphis
me $\underline{J}\hand{\to}\underline{O}_X$ se reconstruit en effet en appliquant le foncteur $i_{\hand{\unsure{*|\relbar}}}$ à
$\underline{J}\,\hand{|}\,X-Z" \hand{\unsure{\longrightarrow|\relbar}} \underline{O}_X\,\hand{|}\,X-Z"$ . Or comme $\underline{J}$ est inversible, il est plat sur S \unsure{e|a}t
 $x\hand{\in} Z" \hand{\Rightarrow} \mathrm{prof}_x\,\underline{J}_{s\,\hand{\unsure{*|x}}} \hand{\geqslant} 2$. Il suffit donc de prouver le
\uline{Lemme} \hand{3.2}. Soit $X\hand{\to} S$ loc de présentation finie, F un Module de présentation
finie sur X plat rel à S, T une partie fermée de X, supposons que pour
\placedmargin{cf\ \S 20.11}
tout $x\hand{\in} X$ au dessus de $s\hand{\in} S$, on ait \hand{\unsure{/}} $\mathrm{prof}_x F_{s\,}\hand{\geqslant} 1$ (resp. $\mathrm{prof}_x F_s \hand{\geqslant} 2$).
Alors l'homomorphisme canonique $F\hand{\to} i_{\hand{*}}(F\hand{|}X-T)$ est injectif (resp.
bijectif), où $i:X-T\hand{\to} T$ est l'immersion canonique.

Démonstration du lemme: on peut supposer S,X affines, et on se
ramène par le procédé bréveté au cas S noethérien. Alors l'hypothèse
implique, en vertu du par.6 , que l'on a $\mathrm{prof}_x F \hand{\geqslant} \mathrm{prof}_x F_s$ pour tout .
$x\hand{\in} X$ sur $s\hand{\in} S$, donc $\mathrm{prof}_x F_s \hand{\geqslant} 1$ (resp.$\hand{\geqslant} 2$) si $x\hand{\in} T$. On conclut alors
grâce au par.5 . (NB Pour bien faire, ce lemme devrait figurer au par.
11 sous la rubrique; élimination des hypothèses noethériennes ...).

Il reste enfin à vérifier la dernière assertion de prop.2, savoir
que X-Z est exactement le domain\unsure{|e} de définition de l"application rationel-%
le rel à S définie par f. Soit $U\hand{\supset} X-Z$ ce domaine, il résulte alors de
la proposition 1 que $U\hand{\longrightarrow} P$ est associé à unsystème lin. de diviseurs
D' sur U\unsure{/|x}S paramétré par $\overset{\hand{\vee}}{P}$ , et on a $D'\hand{|}(X-Z)\times_S\overset{\hand{\vee}}{P} = D\hand{|}(X-Z)\times_S\overset{\hand{\vee}}{P}$ . Appli
quant le résultat d'unicité \hand{(}déjà démontré\hand{)}à D' et $D\hand{|}U\times_S\overset{\hand{\vee}}{P}$ , on voit donc
que ces deux derniers diviseurs sont égaux, donc $D\hand{|}U\times_S\overset{\hand{\vee}}{P}$ n'a pas de
point fixes i.e. $Z\hand{\cap} U=\hand{\emptyset}$ donc U=X-Z , cqfd.

\mspage{\hand{118}}
